\subsection{Проблемы, выявленные и решённые в ходе работы}
\label{subsec:issues}

В ходе разработки модуля, развёртывания окружения и проведения экспериментов
выявлен ряд проблем; их систематизация полезна как практический результат работы.
Проблемы сгруппированы по этапам.

\subsubsection{Сборка и локальное тестирование}

\begin{enumerate}
    \item \textbf{Совместимость Java~17, Flink и DataStream API на клиентском
        процессе.} Запуск задания падал с ошибкой доступа к закрытому полю модуля
        \texttt{java.base} при инициализации транслятора приёмника DataStream.
        Причина --- параметры \texttt{--add-opens} были заданы только для менеджеров
        задач, но не для клиентского процесса. Решение --- задать
        \texttt{env.java.opts.all} (применяется ко всем JVM: менеджер заданий,
        менеджеры задач, клиент).
    \item \textbf{Идемпотентность повторных запусков.} Повторный запуск поверх
        существующих файлов наращивал LSM-дерево и искажал метрики. Решение ---
        принудительная очистка директории таблицы в хранилище перед каждым запуском.
    \item \textbf{Память вне кучи JVM для нативных буферов Vortex.} Чтение Vortex с
        конфигурацией памяти по умолчанию приводило к ошибке исчерпания прямой
        (off-heap) памяти при выделении буферов Arrow/JNI; Parquet и ORC на той же
        конфигурации проходили. Вывод --- Vortex принципиально требует увеличенного
        бюджета off-heap-памяти; конфигурацию менеджеров задач необходимо настраивать
        с учётом нативной стороны.
    \item \textbf{Кросс-компиляция нативной библиотеки.} Сборка нативной библиотеки
        Vortex требовала более новой версии системной библиотеки C, чем доступна на
        целевой версии операционной системы кластера; решение --- кросс-компиляция
        под более старую версию (конфигурация зафиксирована в сборочных контейнерах).
    \item \textbf{Целевая файловая система.} Установлено, что нативная библиотека
        Vortex рассчитана на объектное хранилище, совместимое с S3; распределённая
        файловая система Hadoop (HDFS) и схема URI \texttt{hdfs://} не поддерживаются.
        Это определило выбор бэкенда хранения.
\end{enumerate}

\subsubsection{Развёртывание кластера и автоматизация}

\begin{enumerate}
    \item \textbf{Гонка между завершением быстрого задания и его проверкой.} Плейбук
        проверки ожидал появления задания в состоянии «выполняется», но быстрые
        пакетные задания успевали завершиться до первого опроса. Решение --- извлекать
        идентификатор задания из вывода клиента и опрашивать его конечное состояние,
        принимая в том числе «завершено», и отдельно сигнализировать только об ошибке.
    \item \textbf{Отсутствие классов Hadoop при десериализации графа задания.}
        Перезапуск кластера через обновление конфигурации не экспортировал путь к
        библиотекам Hadoop, в отличие от обычного запуска; решение --- добавить
        экспорт переменных окружения в задачи перезапуска менеджеров.
    \item \textbf{Самоуничтожение shell-обёрток при остановке процессов.} Команды
        вида \texttt{pkill -f 'TaskManagerRunner'} матчили собственную командную
        строку shell-обёртки и посылали сигнал самим себе. Решение --- сузить шаблон
        до полного имени класса и применить приём с символьным классом в регулярном
        выражении (\texttt{[T]askManagerRunner}), который соответствует процессу, но
        не буквальному тексту в командной строке обёртки.
    \item \textbf{Kerberos для доступа к хранилищу.} Для личной учётной записи на
        стенде недоступен режим keytab; пришлось использовать режим кэша билетов с
        отдельным плейбуком их продления, выполняемым первым во всех остальных
        плейбуках.
    \item \textbf{Запись файла метрик на узле менеджера заданий.} Драйвер задания
        выполняется на узле менеджера заданий, где не существует локально созданной
        директории для метрик; решение --- создавать родительский каталог
        непосредственно из кода задания.
    \item \textbf{Автоназначение портов для двух менеджеров задач на узле.} Два
        менеджера задач на одном узле регистрируются только при незаданном порте RPC
        (автоназначение); явное задание порта приводит к конфликту.
\end{enumerate}

\subsubsection{Хранилище: каталог Apache Paimon поверх Apache Ozone}

При повторном создании каталога возникала ошибка «путь должен быть директорией»:
объектное хранилище хранит «пустые директории» как нулевой объект с завершающим
слешем, и при последующем запросе слой совместимости с S3 принимал этот объект-маркер
за обычный файл. Решение --- удалять именно объект-маркер, не затрагивая дочерние
префиксы, после чего обращение к пути корректно разрешается как директория по
дочерним объектам; в перспективе --- использовать собственную схему URI хранилища или
размещать склад в корне бакета.

\subsubsection{Настройка потокового сценария}

\begin{enumerate}
    \item \textbf{Залипание генератора данных на параллелизме~1.} Без указания
        ограничения скорости генератор не шардировался по подзадачам, и пропускная
        способность залипала на пределе одного потока. Решение --- при отсутствии
        ограничения подставлять заведомо большой лимит, что активирует шардирование.
    \item \textbf{Потеря файлов под читателем при частых фиксациях.} Политики
        хранения снимков и манифестов по умолчанию при коротком интервале чекпойнтов
        агрессивно удаляли старые снимки; стартующий писатель или параллельный
        читатель обращались к уже удалённому манифесту. Решение --- расширенные
        политики хранения и асинхронный режим очистки.
    \item \textbf{Исчерпание прямой памяти при высоком параллелизме чтения.}
        Конфигурация с одним крупным менеджером задач на узел упиралась в лимит
        прямой памяти при большом параллелизме чтения и параллельной записи; решение
        --- два менеджера задач на узел меньшего размера (см.
        подраздел~\ref{subsec:deployment}).
    \item \textbf{Накладные расходы JNI при чтении Vortex.} При размере пакета по
        умолчанию (\num{1024} строк) полное сканирование бакета порождало тысячи
        пересечений границы JNI на файл. Решение --- увеличить размер пакета записи
        (он же определяет размер возвращаемого при чтении пакета) до десятков тысяч
        строк, что сократило число пересечений границы в десятки раз и устранило
        отставание Vortex по «чистому» сканированию.
    \item \textbf{Влияние принудительного полного слияния перед итоговыми запросами.}
        Попытка выполнять полное слияние перед фазой FINAL приводила к тому, что
        запросы стартовали до фактического завершения слияния (системная процедура
        слияния запускается как отдельное задание движка, а ожидание её результата
        возвращалось преждевременно); кроме того, агрессивное слияние конкурировало с
        читателем за ввод-вывод. От принудительного слияния отказались: фаза FINAL
        читает состояние «как есть после остановки записи», что и соответствует
        реальной витрине.
\end{enumerate}

\subsubsection{Отклонённые гипотезы по тонкой настройке}

\begin{enumerate}
    \item \textbf{Строчный формат для свежих файлов уровня~0.} Гипотеза «дешёвая
        построчная запись свежих дельт, колоночный формат --- только на нижних
        уровнях» дала прирост пропускной способности записи, но замедлила
        аналитические запросы Vortex на итоговом снимке (итератор слияния вынужден
        совмещать построчные дельты с колоночными файлами, нативный конвейер не
        активируется); при достаточном параллелизме кластера в этом приёме нет нужды.
    \item \textbf{Агрессивные настройки размера файлов и страниц.} Гипотеза
        «крупнее файлы --- лучше амортизация накладных расходов» подтвердилась для
        Parquet (его векторизованный читатель оптимален на крупных страницах), но не
        для Vortex; преимущества Vortex на ряде запросов при этом исчезли, а объём
        хранилища вырос из-за накопления удаляемых файлов в окне хранения при более
        частом слиянии. Использованы значения по умолчанию.
    \item \textbf{Агрессивное снижение порога слияния.} Снижение порога запуска
        слияния учащает фоновые слияния, что увеличивает конкуренцию за ввод-вывод с
        читателем в фазе DURING (наблюдались единичные многосекундные провалы времени
        запросов) и раздувает окно хранения снимков. Использованы значения по умолчанию
        (порог слияния, целевой размер файла).
\end{enumerate}

Что касается записи на наборе TPC-DS --- провести её измерение не удалось: загрузка
крупных наборов в хранилище упиралась в его пропускную способность, что не позволяло
корректно сопоставить форматы; этот аспект покрывается потоковым сценарием.
